\chapter{Hardware Primitive Accelerators}
\label{chap:hardware_accelerators}

\section{Overview of Dedicated Acceleration Engines}
The core computational strength of spGPU resides in its suite of specialized, integer-arithmetic hardware accelerators instantiated within each \texttt{spCORE}. Rather than relying on a general-purpose ALU to execute iterative software rasterization loops, these units are synthesized directly as dedicated finite state machines and data-path pipelines.

Figure~\ref{fig:accelerators} presents the internal mathematical structures and operational data paths of the key acceleration engines.

\begin{figure}[htbp]
\centering
\resizebox{0.95\textwidth}{!}{
    \begin{tikzpicture}[
    >=Stealth,
    scale=0.9, every node/.style={transform shape},
    every text node part/.style={align=center},
    box/.style={draw=black!80, thick, rounded corners=3pt, fill=white, inner sep=5pt, font=\small},
    accbox/.style={box, fill=amber!10, draw=amber!80},
    mathnode/.style={box, fill=cyan!10, draw=cyan!80, font=\scriptsize},
    bus/.style={draw=black!70, line width=1.2pt, ->}
]

% ==================== SUBFIGURE A: BRESENHAM LINE ====================
\node[accbox, font=\bfseries, minimum width=6.5cm, fill=blue!10, draw=blue!70] (title_line) at (3.25, 6.5) {
    (a) Line Accelerator (\texttt{line\_acc.vhd})
};

\node[mathnode] (line_init) at (1.5, 5.2) {
    \textbf{Initialization}\\
    $\Delta X = |X_2 - X_1|$\\
    $\Delta Y = |Y_2 - Y_1|$\\
    $\text{err} = \Delta X - \Delta Y$
};

\node[mathnode] (line_iter) at (5.0, 5.2) {
    \textbf{Single-Cycle Step}\\
    $e_2 = 2 \cdot \text{err} = (\text{err} \ll 1)$\\
    If $e_2 > -\Delta Y$: $\text{err} \gets \text{err} - \Delta Y$, $X \gets X + S_x$\\
    If $e_2 < \Delta X$: $\text{err} \gets \text{err} + \Delta X$, $Y \gets Y + S_y$
};

\node[box, fill=emerald!15, draw=emerald!80] (line_out) at (3.25, 3.8) {
    \textbf{Pixel Output}: $(X, Y, \text{Color}, \text{Valid})$ (1 Pixel / Clock Cycle)
};

\draw[bus] (line_init) -- (line_iter);
\draw[bus] (line_iter) -- (line_out);

% ==================== SUBFIGURE B: TRIANGLE EDGE CROSS PRODUCTS ====================
\node[accbox, font=\bfseries, minimum width=7.0cm, fill=emerald!10, draw=emerald!70] (title_tri) at (11.5, 6.5) {
    (b) Triangle Rasterizer (\texttt{edge\_fill\_top\_v2.vhd})
};

\node[mathnode] (tri_bbox) at (11.5, 5.4) {
    \textbf{Bounding Box Evaluation (\texttt{minmax.vhd})}\\
    $X_{\min} = \min(X_1, X_2, X_3), \quad X_{\max} = \max(X_1, X_2, X_3)$\\
    $Y_{\min} = \min(Y_1, Y_2, Y_3), \quad Y_{\max} = \max(Y_1, Y_2, Y_3)$
};

\node[mathnode, fill=orange!10, draw=orange!80] (tri_cross) at (11.5, 3.8) {
    \textbf{Parallel 2D Cross-Product Units (\texttt{edge\_fill\_v2.vhd})}\\
    $W_0 = (P_x - X_2)(Y_3 - Y_2) - (P_y - Y_2)(X_3 - X_2)$\\
    $W_1 = (P_x - X_3)(Y_1 - Y_3) - (P_y - Y_3)(X_1 - X_3)$\\
    $W_2 = (P_x - X_1)(Y_2 - Y_1) - (P_y - Y_1)(X_2 - X_1)$
};

\node[box, fill=teal!15, draw=teal!80] (tri_inside) at (11.5, 2.2) {
    \textbf{Half-Space Sign Evaluation}\\
    If $\text{sign}(W_0) = \text{sign}(W_1) = \text{sign}(W_2) = \text{sign}(\text{Area})$: Pixel Inside!
};

\draw[bus] (tri_bbox) -- (tri_cross);
\draw[bus] (tri_cross) -- (tri_inside);

% ==================== SUBFIGURE C: CIRCLE SCANLINES ====================
\node[accbox, font=\bfseries, minimum width=6.5cm, fill=purple!10, draw=purple!70] (title_circ) at (3.25, 2.6) {
    (c) Filled Circle Accelerator (\texttt{filled\_circle\_acc.vhd})
};

\node[mathnode] (circ_desc) at (3.25, 1.2) {
    \textbf{Horizontal Scanline Span Filling}\\
    Maintains Midpoint Decision Variable $d \gets d + 4(X-Y) + 14$\\
    Iterates horizontal spans between symmetric octants:\\
    Row $Y_c \pm Y$: Span from $X_c - X$ to $X_c + X$\\
    Row $Y_c \pm X$: Span from $X_c - Y$ to $X_c + Y$
};

\end{tikzpicture}

}
\caption{Mathematical Data Paths of the Hardware Accelerators: (a) Bresenham Line Unit, (b) Triangle Edge Equations Rasterizer, and (c) Filled Circle Scanline Generator.}
\label{fig:accelerators}
\end{figure}

\section{Direct Pixel Writer (\texttt{DRAWPIXEL})}
The \texttt{DRAWPIXEL} operation is executed as a single-cycle direct register latch. When \texttt{acc\_enable\_vec(0)} is asserted, the execution unit registers $X_1$ and $Y_1$ directly to the output buses, asserts \texttt{pixel\_valid\_wire(0) <= '1'}, and pulses \texttt{acc\_finish\_vec(0) <= '1'} in the very next clock cycle. This yields a single-cycle write latency.

\section{Bresenham Line Accelerator (\texttt{line\_acc.vhd})}
The line drawing engine implements an incremental, integer-only version of \textbf{Bresenham's Line Algorithm}. 

\subsection{Mathematical Formulation and Data Path}
Given two endpoints $(X_1, Y_1)$ and $(X_2, Y_2)$, the hardware computes:
\begin{equation}
\Delta X = |X_2 - X_1|, \quad \Delta Y = |Y_2 - Y_1|
\end{equation}
\begin{equation}
S_x = \begin{cases} +1 & \text{if } X_1 < X_2 \\ -1 & \text{otherwise} \end{cases}, \quad S_y = \begin{cases} +1 & \text{if } Y_1 < Y_2 \\ -1 & \text{otherwise} \end{cases}
\end{equation}
The decision error variable is initialized as $\text{err}_0 = \Delta X - \Delta Y$.

During the active \texttt{DRAW} state, the unit evaluates $e_2 = 2 \cdot \text{err}$ (computed via a zero-cost 1-bit left shift) and updates the coordinates and error accumulator in a single clock cycle:
\begin{align}
\text{If } e_2 > -\Delta Y: & \quad \text{err} \gets \text{err} - \Delta Y, \quad X \gets X + S_x \\
\text{If } e_2 < \Delta X: & \quad \text{err} \gets \text{err} + \Delta X, \quad Y \gets Y + S_y
\end{align}
The accelerator outputs exactly \textbf{one pixel per clock cycle} along the line path until $(X, Y) = (X_2, Y_2)$, at which point it transitions to \texttt{DONE} and asserts \texttt{finish <= '1'}.

\section{Midpoint Circle Accelerator (\texttt{circle\_acc.vhd})}
The hollow circle engine implements the \textbf{Midpoint Circle Algorithm}, exploiting the 8-fold rotational and reflective symmetry of the circle.

\subsection{Eight-Octant Traversal}
Given center $(X_c, Y_c)$ and radius $R$, the decision variable is initialized to $d = 3 - 2R$. For each radial step along the first octant, the FSM traverses states \texttt{DRAW1} through \texttt{DRAW8}, emitting the eight symmetric points across eight consecutive clock cycles:
\begin{align}
(X_c + X, Y_c + Y), \quad (X_c - X, Y_c + Y), \quad (X_c + X, Y_c - Y), \quad (X_c - X, Y_c - Y) \\
(X_c + Y, Y_c + X), \quad (X_c - Y, Y_c + X), \quad (X_c + Y, Y_c - X), \quad (X_c - Y, Y_c - X)
\end{align}
The decision parameter $d$ is subsequently updated using addition and shift operations:
\begin{equation}
d_{k+1} = \begin{cases}
d_k + 4(X - Y) + 10, \quad Y \gets Y - 1 & \text{if } d_k > 0 \\
d_k + 4X + 6 & \text{if } d_k \le 0
\end{cases}
\end{equation}
The cycle repeats until $X > Y$.

\section{Filled Circle Accelerator (\texttt{filled\_circle\_acc.vhd})}
To render solid circular discs, \texttt{F\_CIRCLE\_ACC} extends the midpoint circle algorithm into a \textbf{Horizontal Span Rasterizer}.

Instead of plotting isolated boundary pixels, the engine utilizes horizontal scanlines that sweep continuously between opposing symmetric octants:
\begin{itemize}
    \item \textbf{Vertical Octant Pair}: Sweeps horizontal index $i$ from $(X_c - X)$ to $(X_c + X)$ across the two rows $Y_c + Y$ and $Y_c - Y$.
    \item \textbf{Horizontal Octant Pair}: Sweeps horizontal index $j$ from $(X_c - Y)$ to $(X_c + Y)$ across the two rows $Y_c + X$ and $Y_c - X$.
\end{itemize}
This dual-span interpolation guarantees complete geometric coverage with zero interior voids, emitting one filled pixel per clock cycle across each horizontal span.

\section{Filled Triangle Rasterizer (\texttt{edge\_fill\_top\_v2.vhd})}
The triangle fill accelerator utilizes the \textbf{Half-Space Edge Equations Method} (equivalent to 2D cross-products or barycentric orientation testing).

\subsection{Bounding Box Determination (\texttt{minmax.vhd})}
Before rasterizing, the combinational module \texttt{minmax.vhd} evaluates the 2D bounding box of the three triangle vertices:
\begin{equation}
X_{\min} = \min(X_1, X_2, X_3), \quad X_{\max} = \max(X_1, X_2, X_3)
\end{equation}
\begin{equation}
Y_{\min} = \min(Y_1, Y_2, Y_3), \quad Y_{\max} = \max(Y_1, Y_2, Y_3)
\end{equation}
Rasterization scanning is strictly confined within this bounding rectangle $[X_{\min}, X_{\max}] \times [Y_{\min}, Y_{\max}]$, avoiding unnecessary computation across the rest of the canvas.

\subsection{Parallel Edge Function Evaluation (\texttt{edge\_fill\_v2.vhd})}
For any point $P = (X, Y)$, the oriented edge function $E(P)$ of the directed line segment $(X_a, Y_a) \to (X_b, Y_b)$ is defined as:
\begin{equation}
E(P) = (P_x - X_a)(Y_b - Y_a) - (P_y - Y_a)(X_b - X_a)
\end{equation}

Three parallel instances of \texttt{edge\_fill\_v2} compute the edge functions $W_0, W_1, W_2$ for the three edges of the triangle simultaneously:
\begin{align}
W_0 &= (X - X_2)(Y_3 - Y_2) - (Y - Y_2)(X_3 - X_2) \\
W_1 &= (X - X_3)(Y_1 - Y_3) - (Y - Y_3)(X_1 - X_3) \\
W_2 &= (X - X_1)(Y_2 - Y_1) - (Y - Y_1)(X_2 - X_1)
\end{align}
A fourth instance computes the reference oriented area of the full triangle:
\begin{equation}
\text{Area} = (X_3 - X_1)(Y_2 - Y_1) - (Y_3 - Y_1)(X_2 - X_1)
\end{equation}

\subsection{Half-Space Inclusion Test}
A candidate point $P(X, Y)$ lies strictly inside the triangle if and only if all three edge functions have the same sign as the total oriented area:
\begin{equation}
P \in \triangle \iff \left(\text{sign}(W_0) = \text{sign}(\text{Area})\right) \land \left(\text{sign}(W_1) = \text{sign}(\text{Area})\right) \land \left(\text{sign}(W_2) = \text{sign}(\text{Area})\right)
\end{equation}
When this condition evaluates true, \texttt{pixel\_valid <= '1'} is asserted and the pixel is written to VRAM. The rasterizer then advances $X \gets X + 1$; upon reaching $X_{\max}$, it resets $X \gets X_{\min}$ and increments $Y \gets Y + 1$ until the entire bounding box has been processed.
